% Options for packages loaded elsewhere
\PassOptionsToPackage{unicode}{hyperref}
\PassOptionsToPackage{hyphens}{url}
%
\documentclass[
]{article}
\usepackage{amsmath,amssymb}
\usepackage{lmodern}
\usepackage{iftex}
\ifPDFTeX
  \usepackage[T1]{fontenc}
  \usepackage[utf8]{inputenc}
  \usepackage{textcomp} % provide euro and other symbols
\else % if luatex or xetex
  \usepackage{unicode-math}
  \defaultfontfeatures{Scale=MatchLowercase}
  \defaultfontfeatures[\rmfamily]{Ligatures=TeX,Scale=1}
\fi
% Use upquote if available, for straight quotes in verbatim environments
\IfFileExists{upquote.sty}{\usepackage{upquote}}{}
\IfFileExists{microtype.sty}{% use microtype if available
  \usepackage[]{microtype}
  \UseMicrotypeSet[protrusion]{basicmath} % disable protrusion for tt fonts
}{}
\makeatletter
\@ifundefined{KOMAClassName}{% if non-KOMA class
  \IfFileExists{parskip.sty}{%
    \usepackage{parskip}
  }{% else
    \setlength{\parindent}{0pt}
    \setlength{\parskip}{6pt plus 2pt minus 1pt}}
}{% if KOMA class
  \KOMAoptions{parskip=half}}
\makeatother
\usepackage{xcolor}
\IfFileExists{xurl.sty}{\usepackage{xurl}}{} % add URL line breaks if available
\IfFileExists{bookmark.sty}{\usepackage{bookmark}}{\usepackage{hyperref}}
\hypersetup{
  pdftitle={Evaluación Basada en Panel de Investigadores Interdisciplinarios: Protocolo para Agencias Nacionales},
  pdfauthor={A. Rivero and A.I. Scaffold},
  hidelinks,
  pdfcreator={LaTeX via pandoc}}
\urlstyle{same} % disable monospaced font for URLs
\usepackage{longtable,booktabs,array}
\usepackage{calc} % for calculating minipage widths
% Correct order of tables after \paragraph or \subparagraph
\usepackage{etoolbox}
\makeatletter
\patchcmd\longtable{\par}{\if@noskipsec\mbox{}\fi\par}{}{}
\makeatother
% Allow footnotes in longtable head/foot
\IfFileExists{footnotehyper.sty}{\usepackage{footnotehyper}}{\usepackage{footnote}}
\makesavenoteenv{longtable}
\setlength{\emergencystretch}{3em} % prevent overfull lines
\providecommand{\tightlist}{%
  \setlength{\itemsep}{0pt}\setlength{\parskip}{0pt}}
\setcounter{secnumdepth}{-\maxdimen} % remove section numbering
\ifLuaTeX
  \usepackage{selnolig}  % disable illegal ligatures
\fi

\title{Evaluación Basada en Panel de Investigadores Interdisciplinarios:
Protocolo para Agencias Nacionales}
\author{A. Rivero and A.I. Scaffold}
\date{2026}

\begin{document}
\maketitle
\begin{abstract}
Las agencias nacionales de evaluación enfrentan un problema estructural
cuando deben valorar investigadores clasificados como
``Interdisciplinarios'': los paneles monodisciplinares no evalúan de
manera justa trayectorias que abarcan varios campos. Proponemos un
protocolo basado en un panel de tres componentes --- diversidad
disciplinar, coherencia de red y efecto transcampo --- para caracterizar
el tipo y grado de cruce de fronteras disciplinares. Con perfiles
simulados mostramos cómo el panel orienta la composición del comité y
distingue integradores genuinos, polímatas y especialistas mal
clasificados. Identificamos seis modos de fallo y proponemos
salvaguardas.
\end{abstract}

\hypertarget{introducciuxf3n}{%
\section{Introducción}\label{introducciuxf3n}}

Las agencias nacionales evalúan investigadores para promoción,
financiación y acreditación institucional. El procedimiento estándar usa
paneles monodisciplinares, adecuado cuando la cartera del investigador
se ubica en un único campo reconocido.

El problema aparece con investigadores clasificados en el grupo
``Interdisciplinario'' (por ejemplo en esquemas como ANECA o ANVUR):
asignar un panel único crea desajuste estructural entre la pericia del
comité y la trayectoria evaluada.

Este protocolo usa un panel de tres componentes --- diversidad
(\(\Delta\)), coherencia (\(S\)), efecto transcampo (\(E\)) --- para
caracterizar perfiles y diseñar el comité. Operativamente tratamos
``interdisciplinario'' como reclamo de integración, no solo de amplitud:
diversidad alta puede reflejar integración o yuxtaposición
multidisciplinaria desconectada.

\hypertarget{el-panel-indicador}{%
\section{El Panel Indicador}\label{el-panel-indicador}}

\textbf{Diversidad} (\(\Delta\)): índice Rao-Stirling sobre referencias
citadas.

\textbf{Coherencia} (\(S\)): fuerza media de acoplamiento bibliográfico
entre publicaciones del investigador.

\textbf{Efecto transcampo} (\(E\)): fracción de citas recibidas fuera
del campo principal.

La idea central es que ningún componente por sí solo basta. Solo el
triple (\(\Delta\), \(S\), \(E\)) separa integrador, polímata y
especialista.

\hypertarget{protocolo-de-evaluaciuxf3n}{%
\section{Protocolo de Evaluación}\label{protocolo-de-evaluaciuxf3n}}

\hypertarget{paso-1-calcular-el-panel}{%
\subsection{Paso 1: Calcular el Panel}\label{paso-1-calcular-el-panel}}

Para cada investigador en la categoría ``Interdisciplinario'', calcular
(\(\Delta\), \(S\), \(E\)) con datos de publicaciones y citas.
Requisitos mínimos:

\begin{itemize}
\tightlist
\item
  Clasificación disciplinar de referencias.
\item
  Matriz de similitud entre categorías.
\item
  Datos de citación para \(E\).
\end{itemize}

\hypertarget{paso-2-clasificar-el-perfil}{%
\subsection{Paso 2: Clasificar el
Perfil}\label{paso-2-clasificar-el-perfil}}

\begin{longtable}[]{@{}
  >{\raggedright\arraybackslash}p{(\columnwidth - 6\tabcolsep) * \real{0.3721}}
  >{\raggedright\arraybackslash}p{(\columnwidth - 6\tabcolsep) * \real{0.2093}}
  >{\raggedright\arraybackslash}p{(\columnwidth - 6\tabcolsep) * \real{0.2093}}
  >{\raggedright\arraybackslash}p{(\columnwidth - 6\tabcolsep) * \real{0.2093}}@{}}
\toprule
\begin{minipage}[b]{\linewidth}\raggedright
Clasificación
\end{minipage} & \begin{minipage}[b]{\linewidth}\raggedright
\(\Delta\)
\end{minipage} & \begin{minipage}[b]{\linewidth}\raggedright
\(S\)
\end{minipage} & \begin{minipage}[b]{\linewidth}\raggedright
\(E\)
\end{minipage} \\
\midrule
\endhead
Integrador genuino & \(\geq 0.40\) & \(\geq 0.30\) & \(\geq 0.30\) \\
Polímata (no integrativo) & \(\geq 0.40\) & \(< 0.15\) & \(< 0.15\) \\
Especialista (reclasificar) & \(< 0.35\) & cualquiera & cualquiera \\
Provisional (IC superpuesto) & superpone umbrales & superpone umbrales &
superpone umbrales \\
Ambiguo (revisión completa) & resto de combinaciones & & \\
\bottomrule
\end{longtable}

Los umbrales son ilustrativos y deben calibrarse con distribuciones
empíricas por área y etapa de carrera. Si los intervalos de confianza
cruzan varias clases, el caso se trata como provisional.

\hypertarget{paso-3-componer-el-comituxe9}{%
\subsection{Paso 3: Componer el
Comité}\label{paso-3-componer-el-comituxe9}}

Dado un investigador \(r\) con vector de proporciones por categoría
\(p_r\):

\begin{enumerate}
\def\labelenumi{\arabic{enumi}.}
\tightlist
\item
  Definir categorías primarias: \(K_r = \{i : p_{r,i} \geq \tau\}\), con
  \(\tau = 0.15\).
\item
  Incluir al menos un evaluador por cada categoría en \(K_r\).
\item
  Incluir al menos un evaluador con experiencia interdisciplinaria
  demostrada (\(\Delta > 0.40\)).
\item
  Tamaño base del comité: \(|K_r| + 1\).
\item
  Si el caso es provisional por solapamiento de intervalos, añadir un
  evaluador metodológico para revisión de robustez e incertidumbre.
\end{enumerate}

\hypertarget{demostraciuxf3n-con-perfiles-simulados}{%
\subsection{Demostración con Perfiles
Simulados}\label{demostraciuxf3n-con-perfiles-simulados}}

\begin{longtable}[]{@{}lllll@{}}
\toprule
Investigador & \(\Delta\) & \(S\) & \(E\) & Clasificación \\
\midrule
\endhead
Dra. A (integradora) & 0.558 & 0.733 & 0.600 & Integradora genuina \\
Dr.~B (polímata) & 0.562 & 0.000 & 0.063 & Polímata \\
Dra. D (especialista) & 0.288 & 0.881 & 0.211 & Especialista
(reclasificar) \\
\bottomrule
\end{longtable}

Dra. A y Dr.~B muestran casi la misma diversidad, pero con coherencia y
efecto transcampo opuestos. Esa diferencia es precisamente la que evita
clasificaciones erróneas cuando se usa el panel completo.

\hypertarget{modos-de-fallo-y-salvaguardas}{%
\section{Modos de Fallo y
Salvaguardas}\label{modos-de-fallo-y-salvaguardas}}

\hypertarget{f1-premio-a-la-amplitud-sin-profundidad}{%
\subsection{F1: Premio a la amplitud sin
profundidad}\label{f1-premio-a-la-amplitud-sin-profundidad}}

\textbf{Riesgo:} recompensar \(\Delta\) alta ignorando integración e
impacto.

\textbf{Salvaguarda:} exigir simultáneamente umbrales de \(S\) y \(E\)
para ``integrador''.

\hypertarget{f2-penalizaciuxf3n-por-revistas-no-estuxe1ndar}{%
\subsection{F2: Penalización por revistas no
estándar}\label{f2-penalizaciuxf3n-por-revistas-no-estuxe1ndar}}

\textbf{Riesgo:} castigar publicación interdisciplinaria por usar
revistas con métricas distintas a campos hiper-citados.

\textbf{Salvaguarda:} normalización por campo; no comparar factores de
impacto brutos entre disciplinas.

\hypertarget{f3-normas-de-citaciuxf3n-inconmensurables}{%
\subsection{F3: Normas de citación
inconmensurables}\label{f3-normas-de-citaciuxf3n-inconmensurables}}

\textbf{Riesgo:} comparar campos con órdenes de magnitud de citación
distintos.

\textbf{Salvaguarda:} normalizar \(E\) y métricas de citación por línea
base de campo.

\hypertarget{f4-persistencia-de-mala-clasificaciuxf3n}{%
\subsection{F4: Persistencia de mala
clasificación}\label{f4-persistencia-de-mala-clasificaciuxf3n}}

\textbf{Riesgo:} especialistas permanecen en categoría
interdisciplinaria por inercia administrativa.

\textbf{Salvaguarda:} disparador de reclasificación cuando
\(\Delta < 0.35\).

\hypertarget{f5-escasez-de-datos-en-etapas-tempranas}{%
\subsection{F5: Escasez de datos en etapas
tempranas}\label{f5-escasez-de-datos-en-etapas-tempranas}}

\textbf{Riesgo:} con pocas publicaciones, las estimaciones individuales
son inestables (Nakhoda, Whigham y Zwanenburg, 2023).

\textbf{Salvaguarda:} umbral mínimo de publicaciones (por ejemplo,
\(n \geq 15\)), reporte de intervalos de confianza y clasificación
provisional si los intervalos cruzan fronteras de clase.

\hypertarget{f6-juego-estratuxe9gico-por-coautoruxeda}{%
\subsection{F6: Juego estratégico por
coautoría}\label{f6-juego-estratuxe9gico-por-coautoruxeda}}

\textbf{Riesgo:} inflar \(\Delta\) con coautorías oportunistas sin
integración real.

\textbf{Salvaguarda:} ponderar por publicaciones de autoría
correspondiente y contrastar con \(S\).

\hypertarget{discusiuxf3n}{%
\section{Discusión}\label{discusiuxf3n}}

El protocolo complementa, no reemplaza, juicio experto. Su valor
principal es hacer explícita la estructura de evidencia y garantizar que
la composición del comité refleje la cartera real del investigador.

El panel debe interpretarse como evidencia vectorial (\(\Delta\), \(S\),
\(E\)), no como puntaje compuesto único. Esto reduce los incentivos de
manipulación estratégica y mejora la auditabilidad de las decisiones.

\hypertarget{conclusiones}{%
\section{Conclusiones}\label{conclusiones}}

Proponemos un protocolo transparente y auditable para evaluación
nacional de investigadores interdisciplinarios. El panel (\(\Delta\),
\(S\), \(E\)) permite separar perfiles cualitativamente distintos,
ajustar composición de comités y mitigar seis modos de fallo frecuentes.
El siguiente paso natural es calibrar umbrales con distribuciones reales
por área y etapa de carrera.

\hypertarget{referencias}{%
\section{Referencias}\label{referencias}}

\begin{itemize}
\tightlist
\item
  Nakhoda, M., Whigham, P., and Zwanenburg, S. (2023). Quantifying and
  addressing uncertainty in the measurement of interdisciplinarity.
  \emph{Scientometrics}, 128:6107--6127.
\item
  Rafols, I. (2019). S\&T indicators in the wild: contextualization and
  participation for responsible metrics. \emph{Research Evaluation},
  28(1):7--22.
\item
  Rivero, A. (2026). Measuring interdisciplinarity: A multi-component
  indicator panel for research evaluation. Manuscript.
\item
  Xiang, S., Romero, D. M., and Teplitskiy, M. (2025). Evaluating
  interdisciplinary research: Disparate outcomes for topic and knowledge
  base. \emph{Proceedings of the National Academy of Sciences},
  122(16):e2409752122.
\end{itemize}

\end{document}
